% problem_id: S001-cohomology-lemma-left-dual-derived
% source_id: S001 (Stacks Project)
% source_file: cohomology.tex
% statement_locator: cohomology.tex lines 13064--13071
% proof_locator: cohomology.tex lines 13073--13281
% env: lemma
% id_note: label 跨文件重名 ⇒ id 用文件限定形式
% pairing: tight (verified)
% status: complete (statement + author proof verbatim + reason + locators)
%
\begin{lemma}
\label{lemma-left-dual-derived}
Let $(X, \mathcal{O}_X)$ be a ringed space. Let $M$ be an object
of $D(\mathcal{O}_X)$. If $M$ has a left dual in the monoidal category
$D(\mathcal{O}_X)$ (Categories, Definition \ref{categories-definition-dual})
then $M$ is perfect and the left dual is as constructed in
Example \ref{example-dual-derived}.
\end{lemma}

\begin{proof}
Let $x \in X$. It suffices to find an open neighbourhood $U$ of $x$
such that $M$ restricts to a perfect complex over $U$. Hence during the
proof we can (finitely often) replace $X$ by an open neighbourhood of $x$.
Let $N, \eta, \epsilon$ be a left dual.

\medskip\noindent
We are going to use the following argument several times. Choose any
complex $\mathcal{M}^\bullet$
of $\mathcal{O}_X$-modules representing $M$. Choose a K-flat complex
$\mathcal{N}^\bullet$ representing $N$ whose terms are flat
$\mathcal{O}_X$-modules, see Lemma \ref{lemma-K-flat-resolution}.
Consider the map
$$
\eta : \mathcal{O}_X \to
\text{Tot}(\mathcal{M}^\bullet \otimes_{\mathcal{O}_X} \mathcal{N}^\bullet)
$$
After shrinking $X$ we can find an integer $N$ and for
$i = 1, \ldots, N$ integers $n_i \in \mathbf{Z}$ and sections
$f_i$ and $g_i$ of $\mathcal{M}^{n_i}$ and $\mathcal{N}^{-n_i}$
such that
$$
\eta(1) = \sum\nolimits_i f_i \otimes g_i
$$
Let $\mathcal{K}^\bullet \subset \mathcal{M}^\bullet$ be any subcomplex
of $\mathcal{O}_X$-modules containing the sections $f_i$
for $i = 1, \ldots, N$.
Since
$\text{Tot}(\mathcal{K}^\bullet \otimes_{\mathcal{O}_X} \mathcal{N}^\bullet)
\subset
\text{Tot}(\mathcal{M}^\bullet \otimes_{\mathcal{O}_X} \mathcal{N}^\bullet)$
by flatness of the modules $\mathcal{N}^n$, we see that $\eta$ factors through
$$
\tilde \eta :
\mathcal{O}_X \to
\text{Tot}(\mathcal{K}^\bullet \otimes_{\mathcal{O}_X} \mathcal{N}^\bullet)
$$
Denoting $K$ the object of $D(\mathcal{O}_X)$ represented by
$\mathcal{K}^\bullet$ we find a commutative diagram
$$
\xymatrix{
M \ar[rr]_-{\eta \otimes 1} \ar[rrd]_{\tilde \eta \otimes 1} & &
M \otimes^\mathbf{L} N \otimes^\mathbf{L} M
\ar[r]_-{1 \otimes \epsilon} &
M \\
& &
K \otimes^\mathbf{L} N \otimes^\mathbf{L} M
\ar[u] \ar[r]^-{1 \otimes \epsilon} &
K \ar[u]
}
$$
Since the composition of the upper row is the identity on $M$
we conclude that $M$ is a direct summand of $K$ in $D(\mathcal{O}_X)$.

\medskip\noindent
As a first use of the argument above, we can choose the subcomplex
$\mathcal{K}^\bullet = \sigma_{\geq a} \tau_{\leq b}\mathcal{M}^\bullet$
with $a < n_i < b$ for $i = 1, \ldots, N$. Thus $M$ is a direct
summand in $D(\mathcal{O}_X)$ of a bounded complex and we conclude
we may assume $M$ is in $D^b(\mathcal{O}_X)$. (Recall that the process
above involves shrinking $X$.)

\medskip\noindent
Since $M$ is in $D^b(\mathcal{O}_X)$ we may choose
$\mathcal{M}^\bullet$ to be a bounded above complex of
flat modules (by Modules, Lemma \ref{modules-lemma-module-quotient-flat} and
Derived Categories, Lemma \ref{derived-lemma-subcategory-left-resolution}).
Then we can choose $\mathcal{K}^\bullet = \sigma_{\geq a}\mathcal{M}^\bullet$
with $a < n_i$ for $i = 1, \ldots, N$ in the argument above.
Thus we find that we may assume $M$ is a direct summand in
$D(\mathcal{O}_X)$ of a bounded complex of flat modules.
In particular,  $M$ has finite tor amplitude.

\medskip\noindent
Say $M$ has tor amplitude in $[a, b]$. Assuming $M$ is $m$-pseudo-coherent
we are going to show that (after shrinking $X$) we may assume $M$
is $(m - 1)$-pseudo-coherent. This will finish the proof by
Lemma \ref{lemma-perfect-precise} and the fact that
$M$ is $(b + 1)$-pseudo-coherent in any case.
After shrinking $X$ we may assume there exists a strictly perfect
complex $\mathcal{E}^\bullet$ and a map $\alpha : \mathcal{E}^\bullet \to M$
in $D(\mathcal{O}_X)$ such that $H^i(\alpha)$ is an isomorphism for
$i > m$ and surjective for $i = m$. We may and do assume
that $\mathcal{E}^i = 0$ for $i < m$. Choose a distinguished triangle
$$
\mathcal{E}^\bullet \to M \to L \to \mathcal{E}^\bullet[1]
$$
Observe that $H^i(L) = 0$ for $i \geq m$. Thus we may represent
$L$ by a complex $\mathcal{L}^\bullet$ with $\mathcal{L}^i = 0$
for $i \geq m$. The map $L \to \mathcal{E}^\bullet[1]$
is given by a map of complexes
$\mathcal{L}^\bullet \to \mathcal{E}^\bullet[1]$
which is zero in all degrees except in degree $m - 1$
where we obtain a map $\mathcal{L}^{m - 1} \to \mathcal{E}^m$, see
Derived Categories, Lemma \ref{derived-lemma-negative-exts}.
Then $M$ is represented by the complex
$$
\mathcal{M}^\bullet :
\ldots \to
\mathcal{L}^{m - 2} \to
\mathcal{L}^{m - 1} \to
\mathcal{E}^m \to
\mathcal{E}^{m + 1} \to \ldots
$$
Apply the discussion in the second paragraph to this complex to get
sections $f_i$ of $\mathcal{M}^{n_i}$ for $i = 1, \ldots, N$.
For $n < m$ let $\mathcal{K}^n \subset \mathcal{L}^n$
be the $\mathcal{O}_X$-submodule generated by the sections
$f_i$ for $n_i = n$ and $d(f_i)$ for $n_i = n - 1$.
For $n \geq m$ set $\mathcal{K}^n = \mathcal{E}^n$.
Clearly, we have a morphism of
distinguished triangles
$$
\xymatrix{
\mathcal{E}^\bullet \ar[r] &
\mathcal{M}^\bullet \ar[r] &
\mathcal{L}^\bullet \ar[r] &
\mathcal{E}^\bullet[1] \\
\mathcal{E}^\bullet \ar[r] \ar[u] &
\mathcal{K}^\bullet \ar[r] \ar[u] &
\sigma_{\leq m - 1}\mathcal{K}^\bullet \ar[r] \ar[u] &
\mathcal{E}^\bullet[1] \ar[u]
}
$$
where all the morphisms are as indicated above.
Denote $K$ the object of $D(\mathcal{O}_X)$ corresponding to the complex
$\mathcal{K}^\bullet$.
By the arguments in the second paragraph of the proof we obtain
a morphism $s : M \to K$ in $D(\mathcal{O}_X)$ such that the composition
$M \to K \to M$ is the identity on $M$. We don't know that the
diagram
$$
\xymatrix{
\mathcal{E}^\bullet \ar[r] &
\mathcal{K}^\bullet \ar@{=}[r] &
K \\
\mathcal{E}^\bullet \ar[u]^{\text{id}} \ar[r]^i &
\mathcal{M}^\bullet \ar@{=}[r] &
M \ar[u]_s
}
$$
commutes, but we do know it commutes after composing with the
map $K \to M$. By Lemma \ref{lemma-local-actual} after shrinking $X$ we may
assume that $s \circ i$ is given by a map of complexes
$\sigma : \mathcal{E}^\bullet \to \mathcal{K}^\bullet$.
By the same lemma we may assume the composition of $\sigma$
with the inclusion $\mathcal{K}^\bullet \subset \mathcal{M}^\bullet$
is homotopic to zero by some homotopy
$\{h^i : \mathcal{E}^i \to \mathcal{M}^{i - 1}\}$.
Thus, after replacing $\mathcal{K}^{m - 1}$ by
$\mathcal{K}^{m - 1} + \Im(h^m)$ (note that after doing this
it is still the case that $\mathcal{K}^{m - 1}$ is generated
by finitely many global sections), we see that
$\sigma$ itself is homotopic to zero!
This means that we have a commutative solid diagram
$$
\xymatrix{
\mathcal{E}^\bullet \ar[r] &
M \ar[r] &
\mathcal{L}^\bullet \ar[r] &
\mathcal{E}^\bullet[1] \\
\mathcal{E}^\bullet \ar[r] \ar[u] &
K \ar[r] \ar[u] &
\sigma_{\leq m - 1}\mathcal{K}^\bullet \ar[r] \ar[u] &
\mathcal{E}^\bullet[1] \ar[u] \\
\mathcal{E}^\bullet \ar[r] \ar[u] &
M \ar[r] \ar[u]^s &
\mathcal{L}^\bullet \ar[r] \ar@{..>}[u] &
\mathcal{E}^\bullet[1] \ar[u]
}
$$
By the axioms of triangulated categories we obtain a dotted
arrow fitting into the diagram.
Looking at cohomology sheaves in degree $m - 1$ we see that we obtain
$$
\xymatrix{
H^{m - 1}(M) \ar[r] &
H^{m - 1}(\mathcal{L}^\bullet) \ar[r] &
H^m(\mathcal{E}^\bullet) \\
H^{m - 1}(K) \ar[r] \ar[u] &
H^{m - 1}(\sigma_{\leq m - 1}\mathcal{K}^\bullet) \ar[r] \ar[u] &
H^m(\mathcal{E}^\bullet) \ar[u] \\
H^{m - 1}(M) \ar[r] \ar[u] &
H^{m - 1}(\mathcal{L}^\bullet) \ar[r] \ar[u] &
H^m(\mathcal{E}^\bullet) \ar[u]
}
$$
Since the vertical compositions are the identity in both the
left and right column, we conclude the vertical composition
$H^{m - 1}(\mathcal{L}^\bullet) \to
H^{m - 1}(\sigma_{\leq m - 1}\mathcal{K}^\bullet) \to
H^{m - 1}(\mathcal{L}^\bullet)$ in the middle is surjective!
In particular $H^{m - 1}(\sigma_{\leq m - 1}\mathcal{K}^\bullet) \to
H^{m - 1}(\mathcal{L}^\bullet)$ is surjective.
Using the induced map of long exact sequences of cohomology
sheaves from the morphism of triangles above, a diagram chase
shows this implies $H^i(K) \to H^i(M)$ is an isomorphism
for $i \geq m$ and surjective for $i = m - 1$.
By construction we can choose an $r \geq 0$ and a surjection
$\mathcal{O}_X^{\oplus r} \to \mathcal{K}^{m - 1}$. Then the
composition
$$
(\mathcal{O}_X^{\oplus r} \to \mathcal{E}^m \to
\mathcal{E}^{m + 1} \to \ldots ) \longrightarrow
K \longrightarrow M
$$
induces an isomorphism on cohomology sheaves in degrees $\geq m$ and
a surjection in degree $m - 1$ and the proof is complete.
\end{proof}
